日本血吸虫病相关肝纤维化的准确评估是流行地区疾病管理与长期随访的重要基础。超声能够提供非侵入性的影像评估，但复杂的局部回声与解剖结构使精细分级具有挑战。已有深度学习方法能够预测纤维化评分，如何将扫查切面与局部区域线索直接纳入评分过程，并保留可供核查的空间信息，仍需进一步研究。本文提出切面条件区域证据框架 VCRE-Fib，将解剖背景、局部区域信息与整体图像判断结合，用于日本血吸虫病相关肝纤维化的超声精细分级。该框架在空间聚合前形成切面条件化的局部评分证据，通过弱定位引导证据聚合并与全局预测融合，仅凭输入图像即可同时输出纤维化评分、扫查切面及候选异常区域图。我们在经过重新清洗整理、包含 35 个中心共 6,373 名患者、108,709 张超声图像的队列上开发并评估该方法。在按患者划分、包含 4 个中心共 240 名患者、4,107 张图像的测试集上，VCRE-Fib 相对同一数据划分上训练和评估的 SFibAI 将预设综合分级风险降低 7.115\%，图像平均绝对误差由 0.391 降至 0.378，患者最大值聚合、患者中位数聚合及中心平衡风险均有所降低。完整模型的综合分级风险也低于分别移除切面条件化或弱定位的变体。这些结果支持将解剖背景与区域证据纳入超声分级，并在严重程度估计之外提供可供检查的空间预测。
